\newcommand\Fname{XX}
\newcommand\Lname{XX}
\newcommand\Mname{XX}
\newcommand\Group{NN-XX}
\newcommand\Nlist{NN}

\newcommand\Department{XXXX}
\newcommand\Variant{NN}

\usepackage{xstring}

% ця команда отримує позначення дисципліни та обирає приймача роботи.
% я написав сценарій, який визначає це позначення за шляхом до поточного каталогу
% й після ще кількох перевірок генерує шаблон звіту, де заповнене все, крім дати й теми.
% Звісно, не завжди такий шаблонний метод працює, але тоді можна використати \renewcommand
\newcommand{\Work}[1]{
\IfStrEqCase{#1}{
    {<++>}{
        \newcommand\discipline{<++>}
        \IfStrEqCase{\Type}{
            {lab}{
                \newcommand\instructor{<++>} % instructor full name 
                \newcommand\female{1} % gender of instructor
                \newcommand\point{0} % if the report requires point field, declare "point" command. 
                \newcommand\academicRank{XXX} % the report may require instructor academic rank information
                % Optional fields: point, date, topic
            }
        }
    }
}
}
